% Derived from paper3/main.tex
\part{Null-Regime Structural Exclusion}

\section{Compatibility, Instability, and Forced Non-Nullity}

\subsection{Compatibility Operator}
\begin{definition}[Compatibility Operator]\label{def:compatibility-operator}
A map $\Phi: \I \to \E$ is a \textbf{compatibility operator} if it satisfies:
\begin{enumerate}[leftmargin=*]
\item[(C1)] \textbf{Extension Invariance:} For any compatible extension $\mathcal{E}$ of the channel, $\Phi(\I) \cong \Phi(\I_\mathcal{E})$.
\item[(C2)] \textbf{Isomorphism Stability:} If $\I \cong \tilde{\I}$, then $\Phi(\I) \sim \Phi(\tilde{\I})$.
\item[(C3)] \textbf{Lipschitz Lower Bound:} There exists $C > 0$ such that
\[
\dE(\Phi(x), \Phi(\tilde{x})) \geq C \cdot \dw(x,\tilde{x}).
\]
\end{enumerate}
\end{definition}

\begin{definition}[Phenomenological Attractor Set]\label{def:phenomenological-attractor}
For a channel $S$, define
\[
\Attr(S) := \{ e \in \E : e \text{ is stable under all compatible extensions and summable deformations} \}.
\]
\end{definition}

\begin{definition}[Phenomenological Stability]\label{def:phenomenological-stability}
A regime $e \in \E$ is \textbf{stable} for channel $S$ if for every compatible extension $S'$ of $S$ and every summable perturbation family $\{\epsilon_n\}$, the induced regime $\Phi_{S'}(\I')$ stays within a bounded $\dE$-neighborhood of $e$ controlled by $\|\{\epsilon_n\}\|_w$.
\end{definition}

\begin{lemma}[Bridge Lemma]\label{lem:null-bridge}
If $\Phi:\I\to\E$ is a compatibility operator for a CCR channel and $e=\Phi(\I)$ satisfies $\dE(e,\nullphi)=0$, then $e=\nullphi$.
\end{lemma}

\begin{proof}
By the separation property of the phenomenological space, distance zero from $\nullphi$ implies equality with $\nullphi$.
\end{proof}

\subsection{Phenomenological Instability and Forced Non-Nullity}
\begin{theorem}[Phenomenological Instability]\label{thm:phenomenological-instability}
For any CCR channel $S$,
\[
\nullphi \notin \Attr(S).
\]
That is, the null regime is structurally unstable under compatible extensions.
\end{theorem}

\begin{proof}
Assume $\nullphi \in \Attr(S)$ and let $\Phi:\I\to\E$ be a compatibility operator. Because the channel is CCR, finite-energy perturbations preserve the rigid identity object while still allowing admissible deformations in the weighted metric. By the lower-bound clause in Definition~\ref{def:compatibility-operator},
\[
\dE(\Phi(\I),\Phi(\tilde{\I})) \ge C \cdot \dw(\I,\tilde{\I}) > 0
\]
for every non-zero admissible deformation. But if $\nullphi$ were stable, both images would remain equal to $\nullphi$, forcing the left-hand side to be zero. Contradiction.
\end{proof}

\begin{corollary}[Forced Non-Nullity]\label{cor:forced-non-nullity}
Every CCR channel necessarily induces a regime $e_\star \in \E$ with $e_\star \succ \nullphi$.
\end{corollary}

\begin{proof}
Theorem~\ref{thm:phenomenological-instability} rules out $\nullphi$ as a stable admissible attractor. Therefore any admissible stable regime must be non-null.
\end{proof}

\subsection{Positive Regime Structure}
\begin{definition}[Positive Regime Space]\label{def:positive-regime-space}
Define
\[
\Eplus := \{ e \in \E : e \succ \nullphi \}.
\]
\end{definition}

\begin{theorem}[Minimal Positive Regime]\label{thm:minimal-positive-regime}
For every CCR channel there exists a minimal positive regime $e_\star \in \Eplus$.
\end{theorem}

\begin{proof}
Forced non-nullity gives non-emptiness of $\Eplus$. Order-theoretic completeness and the compatibility structure of $\Phi$ yield a minimal element.
\end{proof}

\begin{definition}[Structural Intensity]\label{def:structural-intensity}
For $e \in \Eplus$, define the structural intensity
\[
\iota(e) := \dE(e,\nullphi).
\]
\end{definition}

\begin{theorem}[Universal Intensity Bound]\label{thm:universal-intensity-bound}
If $e_\star$ is the minimal positive regime of a CCR channel, then there exist constants $C,\varepsilon_0>0$ such that
\[
\iota(e_\star) \ge C\varepsilon_0 > 0.
\]
\end{theorem}

\begin{proof}
The lower-bound clause of the compatibility operator transfers any admissible deformation lower bound into positive distance from the null regime. The minimal positive regime therefore has strictly positive structural intensity.
\end{proof}

\begin{theorem}[Forced Phenomenological Closure]\label{thm:forced-phenomenological-closure}
Every CCR system induces a unique minimal phenomenological regime that is non-null, stable, and quantitatively separated from $\nullphi$.
\end{theorem}

\begin{proof}
Combine Corollary~\ref{cor:forced-non-nullity}, Theorem~\ref{thm:minimal-positive-regime}, and Theorem~\ref{thm:universal-intensity-bound}.
\end{proof}

\section{Canonical Families and Structural Closure Consequences}

\begin{proposition}[Profinite Bound]\label{prop:profinite-bound}
In the profinite canonical family, the coupling lower bound induces a positive lower bound on the structural intensity of the minimal positive regime.
\end{proposition}

\begin{proposition}[Symbolic Bound]\label{prop:symbolic-bound}
In the symbolic canonical family, the Bowen-type metric induces the analogous positive lower bound for the minimal positive regime.
\end{proposition}

\begin{theorem}[Rigidity Inheritance]\label{thm:rigidity-inheritance}
If a subsystem or induced family preserves the CCR hypotheses and the compatibility operator, then the inherited regime structure preserves null-regime instability and forced non-nullity.
\end{theorem}

\begin{theorem}[Categorical Universality]\label{thm:null-categorical-universality}
The null-regime exclusion result is universal across admissible CCR realizations in the relevant category of compatible channels.
\end{theorem}

\begin{theorem}[Simulation Lower Bound]\label{thm:null-simulation-lower-bound}
Any computational system attempting to simulate the null-excluding CCR regime structure with arbitrary precision requires unbounded resources as the approximation error tends to zero.
\end{theorem}

\begin{theorem}[Closure by Non-Alternatives]\label{thm:null-closure-by-non-alternatives}
No admissible alternative closes the CCR regime structure while preserving null stability. Any admissible alternative either preserves the same closure structure or exits the CCR class.
\end{theorem}
